\subsection{连续多重线性函数}

定理 1. 设 \(\mathbf{E}_i(\mathbf{i} \leqslant i \leqslant n)\) 与 \(\mathbf{F}\) 是非离散赋值除环 \(\mathbf{K}\) 上的赋范空间，\(f\) 是 \(\prod_{i=1}^{n} \mathbf{E}_i\) 到 \(\mathbf{F}\) 中的多重线性映射。则 \(f\) 在 \(\prod_{i=1}^{n} \mathbf{E}_i\) 上连续，当且仅当存在实数 \(a > 0\) 使得对一切 \(\boldsymbol{x}_i \in \mathbf{E}_i\) （ \(\mathbf{i} \leqslant i \leqslant n\) ）有

\[
\left| \left| f \left(\boldsymbol {x} _ {1}, \boldsymbol {x} _ {2}, \dots , \boldsymbol {x} _ {n}\right) \right| \right| \leqslant a. \left| \left| \boldsymbol {x} _ {1} \right| \right|. \left| \left| \boldsymbol {x} _ {2} \right| \right|. \dots . \left| \left| \boldsymbol {x} _ {n} \right| \right|.\tag{3}
\]

此条件是必要的。因为若 \(f\) 在点 \((\mathbf{o}, \mathbf{o}, \ldots, \mathbf{o})\) 连续，则存在数 \(b > 0\) 使得诸关系 \(||\boldsymbol{x}_i|| \leqslant b\) （ \(\mathbf{i} \leqslant i \leqslant n\) ）蕴含 \(||f(\boldsymbol{x}_1, \ldots, \boldsymbol{x}_n)|| \leqslant \mathbf{i}\) 。设 \(t_0\) 是 \(\mathbf{K}\) 中满足 \(0 < |t_0| < \mathbf{i}\) 的元素；则对每一点 \((\boldsymbol{x}_i) \in \prod_{i=1}^{n} \mathrm{E}_i\) ，只要诸 \(\boldsymbol{x}_i\) 都非零，就存在 \(n\) 个有理整数 \(k_i\) 使 \(b|t_0| < ||t_0^k |\boldsymbol{x}_i|| \leqslant b\) ；于是我们有

\[
\left| t _ {0} \right| ^ {k _ {1} + k _ {2} + \dots + k _ {n}} \left| \left| f \left(\boldsymbol {x} _ {1}, \boldsymbol {x} _ {2}, \dots , \boldsymbol {x} _ {n}\right) \right| \right| \leqslant 1;
\]

另一方面有 \(\frac{\mathbf{I}}{|t_0|^k_i} \leqslant \frac{\mathbf{I}}{b|t_0|} ||\boldsymbol{x}_i||\) ，由此即得关系 (3)，其中 \(a = (\mathbf{I}/b|t_0|)^n\) 。当某个 \(\boldsymbol{x}_i\) 为零时，这个关系显然仍成立。

此条件是充分的。我们将证明：若它成立，则 \(f\) 在每一点 \((\pmb{a}_i)\) （\(\prod_{i=1}^{n} \mathbf{E}_i\) 的）处连续。我们可以写出

\[
f \left(x _ {1}, \dots , x _ {n}\right) - f \left(a _ {1}, \dots , a _ {n}\right) = \sum_ {i = 1} ^ {n} f \left(a _ {1}, \dots , a _ {i - 1}, x _ {i} - a _ {i}, x _ {i + 1}, \dots , x _ {n}\right).
\]

现在，利用 (3)，条件 \(||\pmb{x}_i - \pmb{a}_i|| \leqslant r\) （ \(1 \leqslant i \leqslant n\) ）蕴含

\[
\left| \left| f \left(\boldsymbol {a} _ {1}, \dots , \boldsymbol {a} _ {i - 1}, \boldsymbol {x} _ {i} - \boldsymbol {a} _ {i}, \boldsymbol {x} _ {i + 1}, \dots , \boldsymbol {x} _ {n}\right) \right| \right| \leqslant a r \prod_ {k \neq i} ^ {n} \left(\left| \left| \boldsymbol {a} _ {k} \right| \right| + r\right);
\]

因此，若以 \(c\) 表示诸数 \(\| \pmb{a}_i \| (\mathbf{i} \leqslant i \leqslant n)\) 的最大者，我们有

\[
\left| \left| f \left(\boldsymbol {x} _ {1}, \dots , \boldsymbol {x} _ {n}\right) - f \left(\boldsymbol {a} _ {1}, \dots , \boldsymbol {a} _ {n}\right) \right| \right| \leqslant n a r (c + r) ^ {n - 1}.
\]

由于这个不等式的右端是 r 的常数项为零的多项式，当 r 趋于 o 时它趋于 o；故 f 连续。

注。前面证明过的两个命题是这个定理的推论：双线性函数 \(tx\) 的连续性，依据关系式 \(||tx|| = |t|.||x||\) ；以及命题 7，这只需对 E 的恒等映射应用定理 1，把 E 看作赋有范数 \(p\) 的空间 E 到赋有范数 \(q\) 的空间 E 的线性映射（或反过来）。

